% ============================================================================
% Fichier  : partie4/P04_C15_injections.tex
% Titre    : Injections et Contrôle d'Accès
% Auteur   : MINKA MI NGUIDJOÏ Thierry Emmanuel
% Version  : 1.0 -- novembre 2026
% ============================================================================
% Sources : OWASP Top 10:2021 (A01 Broken Access Control, A03 Injection) ;
% OWASP WSTG. Portail de suivi de LiZbiZ (constat T-05) : cas fictif.
% ============================================================================

\chapter{Injections et Contrôle d'Accès}
\label{chap:injections}
\chaptermark{Injections et contrôle d'accès}

\epigraph{%
    « L'injection confond la donnée et la commande~;\\
    le défaut d'accès confond l'identité et le droit. »%
}{Principe du cours}

\niveauChapitre{Intermédiaire à expert}{Chapitre~\ref{chap:web}~; cours
SEC3266, chapitre sur la sécurité applicative.}

\begin{boxobjectifs}
\begin{itemize}
    \item Expliquer le mécanisme commun à toutes les injections.
    \item Reconnaître l'injection de requête de base de données et l'exécution
          de script dans le navigateur.
    \item Comprendre la référence directe non contrôlée à un objet.
    \item Distinguer contrôle d'accès horizontal et vertical.
    \item Tester sans endommager ni divulguer de données réelles.
    \item Rédiger les constats et leurs corrections.
\end{itemize}
\end{boxobjectifs}

\bigskip

% ============================================================================
\section*{Question directrice : pourquoi changer un seul chiffre dans une adresse donne-t-il accès au colis d'un autre client~?}
% ============================================================================

Sur le portail de suivi, M.~TAMO, connecté avec le compte client de test,
remarque que l'adresse d'une page contient le numéro de son colis. Il
remplace ce numéro par un autre~: le portail affiche le colis d'un autre
client, nom et adresse compris. Aucune injection, aucun outil~: un simple
défaut de contrôle d'accès. C'est le constat T-05.

\medskip
\textbf{Plan du chapitre.}
\begin{enumerate}
    \item Le mécanisme de l'injection (section~\ref{sec:c15-mecanisme}).
    \item Injection de requête et script dans le navigateur
          (section~\ref{sec:c15-familles}).
    \item La référence directe à un objet (section~\ref{sec:c15-idor}).
    \item Horizontal et vertical (section~\ref{sec:c15-acces}).
    \item Tester sans nuire (section~\ref{sec:c15-prudence}).
\end{enumerate}

\separateur

% ============================================================================
\section{Le Mécanisme de l'Injection}
\label{sec:c15-mecanisme}
% ============================================================================

\begin{boxdef}
\termecle{Injection}~: faille par laquelle une donnée fournie par
l'utilisateur est interprétée comme une commande par un système en aval
(base de données, interpréteur, navigateur), faute de séparation entre la
donnée et la commande.
\end{boxdef}

La parade est toujours la même~: séparer la donnée de la commande
(requêtes préparées, échappement selon le contexte, validation stricte des
entrées). Le testeur, lui, cherche les endroits où cette séparation
manque.

% ============================================================================
\section{Injection de Requête et Script dans le Navigateur}
\label{sec:c15-familles}
% ============================================================================

\begin{tableaucours}{L{3.6cm}L{4.4cm}L{4.2cm}}{%
    \tableauHeader{Famille} &
    \tableauHeader{Principe} &
    \tableauHeader{Parade}
}
Injection de requête de base de données & Une entrée modifie la requête
exécutée & Requêtes préparées, paramètres liés \\
Script dans le navigateur & Une entrée renvoyée par la page s'exécute chez
un autre visiteur & Échappement à la sortie, politique de contenu \\
Injection de commande système & Une entrée atteint un interpréteur du
serveur & Éviter l'appel système, listes blanches \\
\end{tableaucours}
\vspace{-0.8em}
\begin{center}{\small\itshape Tableau~15.1. Trois familles d'injection}\end{center}

\begin{boxalert}
\textbf{Violation typique.} Prouver une injection de base de données par
une commande destructrice.

\smallskip
\textbf{Règle~:} on démontre la faille par une preuve inoffensive (délai,
valeur calculée, erreur contrôlée), jamais par une modification ou une
suppression de données. On ne lit pas non plus les données réelles au-delà
de ce qui prouve la faille.
\end{boxalert}

% ============================================================================
\section{La Référence Directe à un Objet}
\label{sec:c15-idor}
% ============================================================================

\begin{boxdef}
\termecle{Référence directe non contrôlée à un objet}~: l'application
désigne un objet (colis, compte, facture) par un identifiant visible, et ne
vérifie pas que l'utilisateur a le droit d'y accéder. Changer l'identifiant
donne accès à l'objet d'autrui.
\end{boxdef}

\begin{boxscenario}{Opération LOGONE. Le colis du voisin}
\textbf{Constatation.} Le portail affiche un colis à partir de son numéro,
passé dans l'adresse, sans vérifier qu'il appartient au client connecté. En
incrémentant le numéro, on accède aux colis d'autres clients~: nom,
téléphone, adresse de livraison.

\medskip
\textbf{Analyse} (confiance élevée). La preuve est faite sur trois colis du
compte de test et un colis voisin, puis arrêtée. L'impact théorique porte
sur l'ensemble des expéditions. C'est le constat T-05, de gravité élevée.
\end{boxscenario}

\begin{boxobjet}{T-05 · Accès indirect aux colis d'autrui}
\textbf{Description~:} le portail ne vérifie pas la propriété du colis
désigné par son numéro. \textbf{Preuve~:} requêtes sur quatre numéros,
captures masquées. \textbf{Gravité~:} élevée. \textbf{Impact~:} divulgation
des données de livraison de tous les clients. \textbf{Recommandation~:}
vérifier, à chaque accès, que l'objet appartient à l'utilisateur
(contrôle côté serveur)~; identifiants non devinables.
\textbf{Détection~:} un compte qui consulte de nombreux colis qui ne sont
pas les siens.
\end{boxobjet}

% ============================================================================
\section{Horizontal et Vertical}
\label{sec:c15-acces}
% ============================================================================

\begin{tableaucours}{L{3.6cm}L{8.8cm}}{%
    \tableauHeader{Défaut d'accès} &
    \tableauHeader{Exemple au portail}
}
Horizontal & Un client accède aux données d'un autre client (T-05) \\
Vertical & Un client accède à une fonction réservée aux agents (mise à jour
d'un colis) \\
\end{tableaucours}
\vspace{-0.8em}
\begin{center}{\small\itshape Tableau~15.2. Deux défauts de contrôle d'accès}\end{center}

% ============================================================================
\section{Tester sans Nuire}
\label{sec:c15-prudence}
% ============================================================================

\begin{boxprocedure}
\textbf{Tester une référence directe à un objet, proprement}
\begin{enumerate}
    \item Repérer l'identifiant d'objet dans l'adresse ou le corps.
    \item Avec un second compte de test, vérifier d'abord sur ses propres
          objets.
    \item Tenter l'accès à un objet voisin~; s'arrêter à la première preuve.
    \item Ne pas parcourir l'ensemble~: la preuve suffit.
    \item Masquer les données réelles dans le rapport.
\end{enumerate}
\end{boxprocedure}

\begin{boxmetier}{Pentesteur et auditeur de sécurité}
\textbf{La preuve minimale.} Une faille d'accès se prouve avec deux comptes
et un objet voisin. L'auditeur qui télécharge toute la base pour « mesurer
l'impact » dépasse son mandat et crée un risque~; l'impact se décrit, il ne
se collecte pas.
\end{boxmetier}

\separateur

\begin{boxresume}
\textbf{Ce qu'il faut retenir de ce chapitre~:}
\begin{itemize}
    \item \textbf{L'injection} confond donnée et commande~; la parade est de
          les séparer.
    \item \textbf{Requête, script, commande}~: trois familles.
    \item \textbf{La référence directe à un objet} donne accès à l'objet
          d'autrui.
    \item \textbf{Horizontal} (autre utilisateur) et \textbf{vertical}
          (fonction réservée).
    \item \textbf{La preuve minimale}~: deux comptes, un objet voisin, puis
          on s'arrête.
\end{itemize}
\end{boxresume}

\begin{boxinfo}
\textbf{Points de contrôle de l'élève.} Je sais~: expliquer le mécanisme
d'injection~; nommer les trois familles~; définir la référence directe à un
objet~; distinguer accès horizontal et vertical~; tester sans nuire.
\end{boxinfo}

\bigskip

\begin{boxexo}[title={\ding{45}\quad Exercice~15.1. Mécanisme \niveaubase}]
En une phrase, qu'ont en commun toutes les injections~?
\end{boxexo}

\begin{boxexo}[title={\ding{45}\quad Exercice~15.2. Familles \niveaubase}]
Associez à chaque famille sa parade principale.
\end{boxexo}

\begin{boxexo}[title={\ding{45}\quad Exercice~15.3. T-05 \niveaubase}]
Pourquoi T-05 n'est-il pas une injection~?
\end{boxexo}

\begin{boxexo}[title={\ding{45}\quad Exercice~15.4. Preuve inoffensive \niveaumoyen}]
Comment prouver une injection de base de données sans modifier ni lire de
données réelles~?
\end{boxexo}

\begin{boxexo}[title={\ding{45}\quad Exercice~15.5. Horizontal ou vertical \niveaumoyen}]
Classez~: un client lit le colis d'un autre~; un client met à jour un colis
alors que seuls les agents le peuvent.
\end{boxexo}

\begin{boxexo}[title={\ding{45}\quad Exercice~15.6. Correction \niveaumoyen}]
Rédigez la correction de T-05 pour l'équipe de développement.
\end{boxexo}

\begin{boxexo}[title={\ding{45}\quad Exercice~15.7. Détection \niveaumoyen}]
Quelle règle de détection signalerait l'exploitation de T-05 en
production~?
\end{boxexo}

\begin{boxexo}[title={\ding{45}\quad Exercice~15.8. Script \niveaumoyen}]
Pourquoi un script renvoyé par la page met-il en danger un autre visiteur,
et non le serveur~?
\end{boxexo}

\begin{boxexo}[title={\ding{45}\quad Exercice~15.9. Procédure \niveauexpert}]
Adaptez la procédure de test de référence directe au cas des factures d'un
portail comptable.
\end{boxexo}

\begin{boxexo}[title={\ding{45}\quad Exercice~15.10. Mandat \niveauexpert}]
Un client demande « la preuve que toutes les données sont exposées ».
Rédigez une réponse qui respecte le principe de la preuve minimale.
\end{boxexo}

\bigskip

\begin{boxinfo}
\textbf{Pour aller plus loin.}
\begin{itemize}
    \item OWASP Top~10~\cite{owasptop10}~; OWASP WSTG~\cite{owaspwstg}.
\end{itemize}
\end{boxinfo}

\begin{boxinfo}
\textbf{Transition.} Les applications modernes exposent des interfaces de
programmation et des applications mobiles~: le Chapitre~\ref{chap:api} leur
est consacré.
\end{boxinfo}
